\section{Independent review and mathematical dialogue}
\label{sec:external-review}

The external reviewer was \textbf{Anthropic's Claude Opus[1m]}
(\artifact{claude-opus-5[1m]}), with \texttt{xhigh} reasoning, running
through \textbf{Claude Code}. The organizers supplied this identification
and configuration. Claude reviewed the manuscript as a separate model
from a different provider, wrote its own checking programs, and discussed
its findings with Codex. We credit Claude with the independent mathematical
review, computational reconstructions, source checks and the suggestion
that led to the strengthened surface example below.

\subsection{What was checked}

In the first report, Claude compared all 46 candidate entries with the
Notebook and read the proofs, verifying a majority in full detail.
It also examined the partial results and prior-work deductions.
Much of the finite checking was reconstructed from the mathematical
descriptions before the authors' programs were available to the reviewer.
The resulting checks therefore provided evidence beyond rerunning the
same implementation. They included finite groups and their characters,
permutation and matrix constructions, polynomial identities, word
algorithms and source-dependent estimates.

Some of the most substantial confirmations concern the carpet results
19.61--19.62. Claude independently rebuilt Chevalley constants and used
its own saturation computation to recover all 104 derived-carpet targets
and all 300 square-completion targets. It also checked the 19 displayed
derivations. For the main examples, it reconstructed the order-2,592
non-monomial group and the order-605 multiple-holomorph example in GAP.
These confirmations support both the displayed arguments and their
finite computational inputs.

After receiving the authors' reply and accessible source material, Claude
conducted a second round. It replayed the documented carpet checks,
verified the order-512 inputs used for 21.114, and checked the finite
cycle data for 15.92. Its final assessment remained acceptance after
revision, with increased confidence in the mathematical content.
The first report found no false theorem among the claims it could settle;
the follow-up resolved several items that had initially lacked sufficient
accessible evidence. The reports give the assessment of each entry.

\subsection{What changed through discussion}

The exchange improved the mathematics as well as its presentation.
Three examples illustrate the different roles of review.

\paragraph{A stronger counterexample: 14.72.}
The original surface had a smooth fixed Cartier divisor, but the reduced
divisor was not globally principal. Claude identified the distinction
and suggested removing a suitable orbit of fibres. Codex worked out
the invariant open \(x^2+1\ne0\), on which
\(f=y-(1+i)x\) generates the fixed ideal while the quotient singularity
survives. Section~\ref{cand:14.72} proves this strengthening and credits
the suggestion. Codex checked the new argument during revision; the exchanged reports
concern the original surface.

\paragraph{An objection corrected: 16.28.}
Claude's first text extraction lost the algebraic-closure bar in
\(\overline{\F_5(t)}\), leading to an objection to the field used in
the example. The authors pointed to the original notation, and Claude
checked the rendered page and withdrew the objection. The argument did
not need repair. The exchange also replaced speculation about proposers'
intentions with explicit convention statements for other entries.

\paragraph{Evidence made accessible.}
Claude confirmed the prime-cycle input to 15.92 once the finite data
were available, and identified the missing small input for the fourth
ancillary carpet checker. The revision supplies that file and explains
the choice of prime. Claude also identified contemporaneous work on
21.106, now acknowledged alongside the proof. These changes improve
what readers can check and how they can assess attribution.

The review substantially increases confidence in the arguments examined.
Its independence concerns the reviewing model and the separately written
checks; it is not a proof of discovery independence, human specialist
endorsement or formal verification. One sizeable verification input,
the \(n=7\) certificate for 20.100, remains outside the public repository,
so the public logs cannot support a fresh complete replay of that case.
We preserve the reports and replies verbatim in
Appendices~\ref{app:review-original}--\ref{app:reply-final}, including
withdrawn assertions. Appendix~\ref{app:review-changes} records the changes.
